\section{Original Whole-Pipeline Cost and Ownership}
\label{sec:cost}
The real FFT's public interface exposes a useful but partial work budget.
A schematic call-cost model is
\begin{equation}
 C_s=C_{\mathrm{input},s}+C_{\mathrm{dispatch},s}
       +b_s C_{\mathrm{butterfly},s}
       +\mathbf{1}_{s=0}C_{\mathrm{prepare}}
       +\mathbf{1}_{s=L-1}C_{\mathrm{reconstruct}}
       +C_{\mathrm{publish},s}.
 \label{eq:cost}
\end{equation}
Here $b_s\le q$ is the number of useful butterflies, and costs need not be
constant across stages or calls. Publication of a previous result may coincide
with preparation of the next. The expression describes locations of work,
not a hardware-independent cycle prediction.

In the source snapshot, real-input packing creates a temporary vector;
copying and bit reversal are bulk operations; reconstruction is bulk; optional
fractional-octave smoothing allocates a prefix-sum array; temporal smoothing
and display-coordinate conversion visit complete arrays. Window-length
reconfiguration may resize storage on the engine path. The existing
implementation therefore does not satisfy an allocation-free, uniformly
small-call contract merely because its inner FFT is resumable.

If $H=\rho N$ for fixed $\rho>0$, then $q=\Theta(\log N)$ for the real FFT.
The butterfly work assigned to ordinary calls grows logarithmically, while
preparation and reconstruction still contribute $\Theta(N)$ operations to
boundary calls. Consequently, the full real-FFT path retains a linear peak
operation count under this scaling. The change from a full
$\Theta(N\log N)$ burst to smaller inner batches can still be useful, but the
asymptotic distinction must be stated at the correct boundary.

\subsection{Magnitude processing and SIMD}
Fractional-octave smoothing uses a prefix sum of magnitudes, allowing a range
mean to be evaluated with two prefix accesses after a linear setup pass. For
fraction $\beta$ of an octave, the nominal band around $f_k$ is
$[f_k2^{-\beta/2},f_k2^{\beta/2}]$, discretized to FFT-bin indices and adjusted
near Nyquist. The result is a visualization-oriented magnitude average, not a
constant-Q transform or a power-conserving spectral estimate. Phase is discarded
when smoothed magnitudes replace complex coefficients.

The Fourier module instantiates arithmetic on \code{simd::float\_4}. Each lane
carries an independent input channel through the same traversal; the SIMD
lanes are not four consecutive bins of one transform. This arrangement
shares index/control work across channels. It does not establish a measured
fourfold speedup, and the scalar experiments in
Appendix~\ref{sec:method} do not validate Rack SIMD instantiations. Rack's plugin API supplies the host processing and SIMD
context~\cite{rack}.

\subsection{Snapshots and publication}
The input ring can continue advancing because \code{buffer()} copies the
selected frame into transform-owned storage. Resuming directly from a mutable
ring without preserving the selected samples would violate
Equation~\eqref{eq:dft}. Any attempt to distribute packing must therefore solve
input lifetime as well as arithmetic scheduling.

Similarly, a display-readiness flag is not by itself an engine-to-display
synchronization protocol. Concurrent reads of arrays being modified require an
ownership and publication design consistent with the C++ memory model. The
original source's shared display buffers are not certified by its scalar tests.
A bounded exchange of preallocated buffers needs an explicit slow-consumer
policy and safe reuse; an atomic pointer alone is insufficient. The successor's
ownership exchange is described in Section~\ref{sec:one-hop}.
